\documentclass{article}

% Language setting
% Replace `english' with e.g. `spanish' to change the document language
\usepackage[english]{babel}

% Set page size and margins
% Replace `letterpaper' with `a4paper' for UK/EU standard size
\usepackage[letterpaper,top=2cm,bottom=2cm,left=3cm,right=3cm,marginparwidth=1.75cm]{geometry}

% Useful packages
\usepackage{amsmath}
\usepackage{graphicx}
\usepackage{subcaption}
\usepackage[colorlinks=true, allcolors=blue]{hyperref}
\usepackage{hyperref}      

\usepackage{caption}
\usepackage{subcaption}

\usepackage{diagbox}
\usepackage{float}
\usepackage{booktabs}
\usepackage{multirow}
\usepackage{tcolorbox}
\tcbset{colback=white, colframe=black, boxrule=0.5pt}

\newcommand{\drug}{\texttt{bio:Drug}}
\newcommand{\disease}{\texttt{bio:Disease}}
\newcommand{\type}{\texttt{rdf:type}}
\newcommand{\domain}{\texttt{rdfs:domain}}
\newcommand{\range}{\texttt{rdfs:range}}
\newcommand{\subclassof}{\texttt{rdfs:subClassOf}}
\newcommand{\affects}{\texttt{bio:affects}}
\newcommand{\plass}{\texttt{rdfs:Class}}
\newcommand{\property}{\texttt{rdf:Property}}

% Set the section counter to start at 0
\setcounter{section}{-1}

% Change section numbering to Alphabetic (A, B, C...)
\renewcommand{\thesection}{\Alph{section}}

% Special rule to show the first section (index 0) as "0"
\renewcommand{\thesection}{\ifnum\value{section}=0 0\else\Alph{section}\fi}

\usepackage{minted}


\usepackage[style=numeric,backend=biber]{biblatex}
\addbibresource{references.bib} % The name of your .bib file

\begin{document}

% --- START OF FULL TITLE PAGE ---
\begin{titlepage}
    \centering
    \vspace*{3cm} % Adjust space from top
    
    {\Huge \textbf{Knowledge Graphs Exploitation} \par}
    \vspace{0.5cm}
    {\LARGE Master's degree in Data Science \\ \LARGE Semantic Data Management \par}
    \vspace{0.3cm}
    {\Large Lab Assignment 3}
    
    \vspace{2.5cm}
    
    {\Large \textbf{Authors:}} \\
    \vspace{0.5cm}
    {\Large Pau Adal Solís \\ Antoni Lopera Barril \par}


    \vspace{3cm}
    
    \includegraphics[width=0.2\textwidth]{UPC logo.png}
    
    \vfill % Pushes the date to the bottom
    
    {\Large \today \par} % Automatically inserts current date
    \vspace{2cm}
\end{titlepage}
% --- END OF FULL TITLE PAGE ---

\tableofcontents




% ====================================================================
\section*{A \hspace{0.5em} Getting Familiar with KGEs}
% ====================================================================

We use the \textbf{PubMed} citation dataset: 19,717 papers (nodes) and 44,338
citation links (edges), with a single relation type \texttt{cites}. The graph is
loaded via PyTorch Geometric's \texttt{Planetoid} and converted into KGE triples
of the form \texttt{(paper\_i, cites, paper\_j)}, then split deterministically
into training (80\%), validation (10\%) and testing (10\%) with
\texttt{random\_state = 2026} for full reproducibility.

% --------------------------------------------------------------------
\subsection*{A.1 \hspace{0.3em} The Most Basic Model}
% --------------------------------------------------------------------

\paragraph{Model and hyperparameters.}
We train a \textbf{TransE} model with PyKEEN. The assignment statement only
requires a \emph{sensible} configuration, so we use the values shown in
Table~\ref{tab:a1-hparams}. These follow the typical defaults recommended in the
PyKEEN documentation and are large enough to obtain meaningful embeddings on a
KG of this size without overfitting.

\begin{table}[H]
\centering
\caption{TransE hyperparameters used in A.1.}
\label{tab:a1-hparams}
\begin{tabular}{ll}
\toprule
\textbf{Hyperparameter} & \textbf{Value} \\
\midrule
Embedding dimension     & 128 \\
Number of training epochs & 100 \\
Batch size              & 512 \\
Random seed             & 2026 \\
\bottomrule
\end{tabular}
\end{table}

\paragraph{Paper selection.}
From the training triples, we filter the papers that have between $8$ and $12$
outgoing citations (this range balances representativeness --- a paper with too
few citations would give a noisy mean --- with generality). We sort the
candidates numerically by paper id and pick the first one. Because the split is
seeded and the sort is deterministic, the chosen paper is always the same in
every execution. We denote this paper as \textbf{A}.

\paragraph{Mathematical reasoning.}
TransE encodes the assumption that
\begin{equation}
    \mathbf{e}(h) + \mathbf{e}(r) \;\approx\; \mathbf{e}(t)
    \qquad \forall (h,r,t) \in \mathcal{KG}.
\end{equation}
Let $B_1,\dots,B_n$ be the papers cited by $A$, i.e., the set of triples
$(A, \texttt{cites}, B_i)$ for $i=1,\dots,n$. We are looking for the most
probable embedding $\mathbf{x}^\star$ of a \emph{hypothetical paper} $X$ that
would cite works similar to those cited by $A$. Applying the TransE
relation translation, each citation $(X, \texttt{cites}, B_i)$ implies
\begin{equation}
    \mathbf{x}^\star + \mathbf{e}(\texttt{cites}) \;\approx\; \mathbf{e}(B_i)
    \;\;\Longleftrightarrow\;\;
    \mathbf{x}^\star \;\approx\; \mathbf{e}(B_i) - \mathbf{e}(\texttt{cites}).
\end{equation}
Since we want one single vector $\mathbf{x}^\star$ that is consistent with all
$n$ constraints simultaneously, we take the mean of the $n$ right-hand sides:
\begin{equation}
    \boxed{\;
    \mathbf{TARGET} \;=\; \mathbf{x}^\star
    \;=\;
    \underbrace{\frac{1}{n}\sum_{i=1}^n \mathbf{e}(B_i)}_{\textstyle \overline{\mathbf{e}_\text{cited}}}
    \;-\; \mathbf{e}(\texttt{cites}).
    \;}
\end{equation}
This is the least-squares optimum: \textbf{TARGET} minimises
$\sum_{i=1}^{n}\Vert (\mathbf{x}+\mathbf{e}(\texttt{cites})) -
\mathbf{e}(B_i)\Vert^{2}$ over $\mathbf{x}$.

\paragraph{Retrieving the closest real paper.}
We then search for the existing paper whose embedding is closest to
\textbf{TARGET} in the learned space. We use \emph{cosine similarity} because
TransE embeddings are unconstrained in magnitude and angular distance is the
standard choice. The candidate $A$ itself is excluded from the search. The
recovered paper is called \textbf{X}:
\begin{equation}
    X \;=\; \arg\!\max_{p \,\in\, \mathcal{V} \setminus \{A\}}
    \cos\!\big(\mathbf{e}(p), \mathbf{TARGET}\big).
\end{equation}

\paragraph{Results.}
Table~\ref{tab:a1-result} summarises the outcome. Figure~\ref{fig:a1-pca} shows
the 2D sketch of the reasoning, obtained by PCA-projecting the embeddings of
$A$, its cited papers $\{B_i\}$, the \textbf{TARGET} vector, and the retrieved
paper $X$.

\begin{table}[H]
\centering
\caption{Result of the closest-paper search in A.1.}
\label{tab:a1-result}
\begin{tabular}{ll}
\toprule
\textbf{Field} & \textbf{Value} \\
\midrule
Chosen paper $A$        & \texttt{<paper\_A>} \\
\# outgoing citations of $A$ & $n \in [8,12]$ \\
Retrieved paper $X$     & \texttt{<paper\_X>} \\
$\cos(\mathbf{e}(X), \mathbf{TARGET})$ & $<$value$>$ \\
\bottomrule
\end{tabular}
\end{table}

\begin{figure}[H]
\centering
\includegraphics[width=0.85\linewidth]{A1_visualization.png}
\caption{2D PCA projection of TransE embeddings. Paper $A$ (blue), papers cited
by $A$ (grey), the computed \textbf{TARGET} vector (orange star) and the
retrieved paper $X$ (red) closest to \textbf{TARGET} by cosine similarity. The
target vector lies in the geometric centre of the cited papers translated by
$-\mathbf{e}(\texttt{cites})$, exactly as predicted by the TransE translation
assumption.}
\label{fig:a1-pca}
\end{figure}

\paragraph{Interpretation.}
The recovered paper $X$ is the paper of the corpus whose embedding most
closely matches the ideal position predicted by TransE for a paper that would
cite works similar to those of $A$. Intuitively, $X$ acts as a
``citation-twin'' of $A$: it is the paper whose citation behaviour the model
expects to be most similar to that of $A$ at the relation level. The companion
code that produces this analysis is provided in the notebook
\texttt{Lab\_Template.ipynb}, section \emph{A.1}.

% --------------------------------------------------------------------
\subsection*{A.2 \hspace{0.3em} Improving TransE}


\paragraph{The drawback: one-to-many relations.}
TransE relies on the strict translation assumption
$\mathbf{e}(h) + \mathbf{e}(r) \approx \mathbf{e}(t)$ \emph{for every} triple.
This works well for one-to-one relations but breaks down for the
\textbf{one-to-many} (and many-to-many) case. In our KG the only relation,
\texttt{cites}, is intrinsically one-to-many: a single paper typically cites
several different papers. Whenever a head $h$ participates in $n>1$ triples
with the same relation $r$,
\begin{equation}
    (h, r, t_1),\, (h, r, t_2),\, \dots,\, (h, r, t_n) \,\in\, \mathcal{KG},
\end{equation}
TransE has to satisfy $n$ constraints simultaneously:
\begin{equation}
    \mathbf{e}(h) + \mathbf{e}(r) \;\approx\; \mathbf{e}(t_1)
    \quad\wedge\quad
    \mathbf{e}(h) + \mathbf{e}(r) \;\approx\; \mathbf{e}(t_2)
    \quad\wedge\;\cdots\;\wedge\quad
    \mathbf{e}(h) + \mathbf{e}(r) \;\approx\; \mathbf{e}(t_n).
\end{equation}
Since the left-hand side is identical in all $n$ equations, the right-hand
sides are forced to coincide:
\begin{equation}
    \mathbf{e}(t_1) \;\approx\; \mathbf{e}(t_2) \;\approx\; \cdots
    \;\approx\; \mathbf{e}(t_n).
\end{equation}
This contradicts the fact that $t_1,\dots,t_n$ are \emph{different} entities
that, in general, participate in many other triples and therefore should be
representable by distinct vectors. The model can only find a global compromise
that increases the residual error and degrades link-prediction quality (low
Hits@1, low MRR, as we confirm empirically in A.3).

\paragraph{Concrete example from our KG.}
Consider any paper $A$ from the PubMed KG (e.g., the paper selected in A.1).
A subset of triples in the training set has the form
\begin{align}
    (A,\, \texttt{cites},\, B_1) \nonumber\\
    (A,\, \texttt{cites},\, B_2) \nonumber\\
    \;\;\vdots\;\;            \nonumber\\
    (A,\, \texttt{cites},\, B_n) \label{eq:one-to-many-example}
\end{align}
with $n$ between 8 and 12 in our selected case. TransE forces
$\mathbf{e}(B_1) \approx \mathbf{e}(B_2) \approx \cdots \approx \mathbf{e}(B_n)$,
even though these papers are heterogeneous: they are cited by many other
papers and each one participates in its own set of triples. The constraint
system has no exact solution and the loss never converges to zero --- this
explains the limited expressive power of TransE on citation-style KGs.

\paragraph{The solution: TransH \cite{wang2014transh}.}
TransH \emph{(Wang et al., AAAI 2014)} fixes the one-to-many problem by
\textbf{projecting} entities onto a \emph{relation-specific hyperplane}
before applying the translation. For each relation $r$, the model learns two
vectors:
\begin{itemize}[leftmargin=1.5em,nosep]
    \item the \textbf{hyperplane normal} $\mathbf{w}_r \in \mathbb{R}^d$ with
          $\Vert\mathbf{w}_r\Vert = 1$;
    \item the \textbf{translation vector} $\mathbf{d}_r \in \mathbb{R}^d$
          that lives \emph{on} the hyperplane.
\end{itemize}
Any entity embedding $\mathbf{e}$ is projected onto the hyperplane of $r$ by
removing its component along the normal:
\begin{equation}
    \mathbf{e}_\perp
    \;=\;
    \mathbf{e} \;-\; \big(\mathbf{w}_r^{\top}\mathbf{e}\big)\,\mathbf{w}_r.
\end{equation}
The triple-plausibility (energy) score becomes
\begin{equation}
    f_r(h, t)
    \;=\;
    \big\Vert\,
        \mathbf{e}_\perp(h) \;+\; \mathbf{d}_r \;-\; \mathbf{e}_\perp(t)
    \,\big\Vert_{2}^{\,2}.
\end{equation}
The translation is now performed on a \emph{relation-dependent} subspace, so
two entities $t_i$ and $t_j$ that share the same head $h$ under relation $r$
can have distinct embeddings in the original space $\mathbb{R}^d$ while their
\emph{projections} on the hyperplane of \texttt{cites} are forced to coincide
with $\mathbf{e}_\perp(h)+\mathbf{d}_r$. In other words, $B_1,\dots,B_n$ can
remain different along the direction $\mathbf{w}_r$ (which is the dimension
\textit{not} used by this relation) and only need to agree on the $d-1$
remaining directions. The contradiction (3) is therefore resolved.

\paragraph{Geometric intuition.}
TransE places all tails $t_i$ at the same point
$\mathbf{e}(h)+\mathbf{e}(r)$, which is impossible to satisfy for several
distinct $t_i$. TransH instead constrains only the projection onto the
hyperplane, leaving one dimension free along $\mathbf{w}_r$ where the tails
can be spread apart. Every relation gets its own hyperplane, so a single
entity can play different roles under different relations without conflicting
constraints. Empirically (see A.3) this richer geometry improves the
modelling of one-to-many citation patterns over plain TransE.


% --------------------------------------------------------------------
\subsection*{A.3 \hspace{0.3em} Training KGEs}

 Five Knowledge Graph Embedding (KGE) models were trained on the PubMed citation
  graph: three translational models (TransE, TransH, RotatE) and two semantic models
  (DistMult, ComplEx). All models were trained using the PyKEEN pipeline with the
  following base configuration:

  \begin{itemize}
      \item \textbf{Embedding dimension:} 128
      \item \textbf{Number of epochs:} 100 (maximum)
      \item \textbf{Batch size:} 512
      \item \textbf{Random seed:} 2026
      \item \textbf{Negative sampling:} Uniform (default)
      \item \textbf{Early stopping:} monitored on validation MRR,
      patience = 5 evaluations, relative delta = 0.002
  \end{itemize}

  Early stopping monitors the Mean Reciprocal Rank (MRR) on the validation set
  after each evaluation. Training is stopped when no improvement greater than 0.2\%
  is observed for 5 consecutive evaluations, and the model weights are restored
  to the best checkpoint found during training. This prevents overfitting and
  reduces unnecessary computation.

  \subsubsection*{Results}

  \begin{table}[h]
  \centering
  \begin{tabular}{lccc}
  \hline
  \textbf{Model} & \textbf{Hits@1} & \textbf{Hits@3} & \textbf{MRR} \\
  \hline
  TransE   & 0.0000 & 0.4141 & 0.2354 \\
  TransH   & 0.0027 & 0.0262 & 0.0296 \\
  RotatE   & \textbf{0.6302} & \textbf{0.7324} & \textbf{0.6952} \\
  DistMult & 0.4579 & 0.5624 & 0.5316 \\
  ComplEx  & 0.0002 & 0.0010 & 0.0071 \\
  \hline
  \end{tabular}
  \caption{Link prediction results on the test set for all trained KGE models.}
  \label{tab:kge_results}
  \end{table}

  \subsubsection*{Analysis}

  \textbf{RotatE} achieves by far the best performance (MRR = 0.6952, Hits@1 = 0.6302).
  This is consistent with its theoretical properties: RotatE models each relation as
  a rotation in complex vector space, which allows it to capture asymmetric and
  directional relations. Since \textit{cites} is inherently asymmetric --- the fact
  that paper A cites paper B does not imply the reverse --- RotatE's geometric
  formulation is well-suited to this KG. Its embeddings are saved as
  \texttt{best\_entity\_embeddings.npy} (shape: 19717 $\times$ 128) for use in
  Section B.5.2.

  \textbf{DistMult} ranks second (MRR = 0.5316). Although DistMult assumes symmetric
  relations through its bilinear score $\text{Score}(h, r, t) = h^\top \text{diag}(r) t$,
  it still captures latent semantic patterns effectively. The rich semantic structure
  of the citation graph compensates for the symmetry assumption to some extent.

  \textbf{TransE} obtains a mediocre MRR of 0.2354 and a Hits@1 of 0.0000, meaning
  it never correctly predicts the exact tail entity. This is a direct consequence of
  the one-to-many limitation discussed in Section A.2: since a paper typically cites
  multiple papers, TransE cannot simultaneously satisfy all translational constraints
  without residual error.

  \textbf{TransH} performs surprisingly worse than TransE (MRR = 0.0296), despite
  being theoretically superior for one-to-many relations. This is likely due to its
  additional parameters (the relation-specific hyperplane vectors $w_r$ and $d_r$),
  which require more training time to converge. With only 100 epochs available and
  early stopping, TransH does not have sufficient training budget to reach its
  optimal configuration.

  \textbf{ComplEx} barely learns (MRR = 0.0071). Like TransH, its complex-valued
  embeddings double the parameter count compared to DistMult, making convergence
  within 100 epochs very difficult for a graph of this size.

  \subsubsection*{Hyperparameter Exploration}

  To understand the effect of key training parameters, RotatE was retrained varying
  one parameter at a time while keeping the rest fixed at the base configuration.
  Early stopping was applied in all runs.

  \paragraph{Embedding Dimension}

  \begin{table}[h]
  \centering
  \begin{tabular}{lccc}
  \hline
  \textbf{Embedding Dim} & \textbf{Hits@1} & \textbf{Hits@3} & \textbf{MRR} \\
  \hline
  64  & 0.6177 & 0.7090 & 0.6783 \\
  128 & 0.6302 & 0.7324 & 0.6952 \\
  \textbf{256} & \textbf{0.6479} & \textbf{0.7382} & \textbf{0.7068} \\
  512 & 0.6258 & 0.7278 & 0.6905 \\
  \hline
  \end{tabular}
  \caption{Effect of embedding dimension on RotatE performance.}
  \end{table}

  As expected, increasing the embedding dimension improves performance up to a point,
  as larger embeddings provide more expressive power to capture complex relational
  patterns. The optimal value is 256, achieving MRR = 0.7068. At dimension 512,
  performance slightly decreases. Since early stopping controls training duration,
  this drop is likely due to the model having too many parameters relative to the
  size of the training set, leading to poorer generalisation despite not overfitting
  in the classical sense.

  \paragraph{Number of Negative Samples per Positive}

  \begin{table}[h]
  \centering
  \begin{tabular}{lccc}
  \hline
  \textbf{Num Negatives} & \textbf{Hits@1} & \textbf{Hits@3} & \textbf{MRR} \\
  \hline
  1  & 0.6302 & 0.7324 & 0.6952 \\
  5  & 0.7055 & 0.7964 & 0.7586 \\
  \textbf{10} & \textbf{0.7364} & \textbf{0.8125} & \textbf{0.7810} \\
  \hline
  \end{tabular}
  \caption{Effect of number of negative samples per positive on RotatE performance.}
  \end{table}

  Increasing the number of negative samples per positive triple consistently and
  significantly improves performance. More negatives provide a stronger and more
  diverse training signal, forcing the model to learn finer distinctions between
  valid and invalid triples. This is both the expected and observed behaviour. The
  improvement from 1 to 10 negatives is substantial (+8.6 MRR points), suggesting
  that the base configuration was under-sampling negatives.

  \paragraph{Loss Margin}

  \begin{table}[h]
  \centering
  \begin{tabular}{lccc}
  \hline
  \textbf{Margin} & \textbf{Hits@1} & \textbf{Hits@3} & \textbf{MRR} \\
  \hline
  \textbf{0.5} & \textbf{0.6578} & \textbf{0.7506} & \textbf{0.7151} \\
  1.0 & 0.6302 & 0.7324 & 0.6952 \\
  2.0 & 0.5408 & 0.6458 & 0.6122 \\
  \hline
  \end{tabular}
  \caption{Effect of loss margin on RotatE performance.}
  \end{table}

  Contrary to the naive expectation that a larger margin enforces stricter separation
  and thus improves performance, smaller margins yield better results. A large margin
  makes the optimisation problem harder: the model is required to produce a very large
  score gap between positive and negative triples, which is difficult to achieve within
  100 epochs on a graph of this complexity. A margin of 0.5 provides sufficient
  separation while remaining achievable, leading to the best MRR of 0.7151.
% --------------------------------------------------------------------
% TODO

% --------------------------------------------------------------------
\subsection*{A.4 \hspace{0.3em} Negative Sampling}
% --------------------------------------------------------------------
% TODO


  KGE models learn to give high scores to positive triples $(h, r, t)$.
  The problem is that if we only show the model positive examples,
  it can cheat by giving a high score to everything without learning
  anything useful. Negative sampling fixes this by creating fake triples
  where either the head or the tail is replaced with a random entity,
  so the model has to learn to tell real triples from fake ones.

  The simplest approach is \textbf{Uniform Sampling} that basically gives the sample probaility of changing the head or tail ($50\%$ chance of
  replacing the head, $50\%$ chance of replacing the tail). It does not
  consider what kind of relation we are dealing with.
  
On the other hand, \textbf{Bernoulli sampling} uses two statistics from the training set to decide
  which side to corrupt more often. The idea is that for a one-to-many
  relation, randomly replacing the tail often produces a plausible triple
  that could actually be true, giving the model an easy negative that
  does not help it learn. The same problem happens on the head side for
  many-to-one relations.
  
 Bernoulli sampling uses two statistics computed from the training set to decide
  which side to corrupt more often:

  \begin{itemize}
    \item $\text{tph}$: average number of tails per head.
    \item $\text{hpt}$: average number of heads per tail.
  \end{itemize}

  \begin{align}
    P(\text{corrupt tail}) &= \frac{\text{tph}}{\text{tph} + \text{hpt}} \\[4pt]
    P(\text{corrupt head}) &= \frac{\text{hpt}}{\text{tph} + \text{hpt}}
  \end{align}

  If the relation is one-to-many ($\text{tph} \gg \text{hpt}$), it corrupts
  the tail more often to generate harder negatives. If it is many-to-one,
  it does the opposite.

  From the PubMed training set (70{,}918 triples, 19{,}717 distinct heads,
  17{,}353 distinct tails):

  \[
    \text{tph} = \frac{70918}{19717} = 3.5968, \qquad
    \text{hpt} = \frac{70918}{17353} = 4.0868
  \]

  \[
    P(\text{corrupt tail}) = \frac{3.5968}{3.5968 + 4.0868} = 0.4681, \qquad
    P(\text{corrupt head}) = \frac{4.0868}{3.5968 + 4.0868} = 0.5319
  \]

  \begin{table}[h]
  \centering
  \begin{tabular}{lcc}
  \hline
  \textbf{Statistic} & \textbf{Uniform} & \textbf{Bernoulli} \\
  \hline
  tph & --- & 3.5968 \\
  hpt & --- & 4.0868 \\
  $P(\text{corrupt head})$ & 0.5000 & 0.5319 \\
  $P(\text{corrupt tail})$ & 0.5000 & 0.4681 \\
  \hline
  \end{tabular}
  \caption{Negative sampling statistics for the \textit{cites} relation.}
  \label{tab:a4}
  \end{table}

  As we can see, the tph and hpt values are quite similar ($3.60$ vs $4.09$),
  so Bernoulli ends up giving almost the same probabilities as uniform
  ($0.53$ vs $0.47$). This makes sense for the \textit{cites} relation because
  in PubMed, papers tend to cite a similar number of papers as they get
  cited by, so the relation is roughly balanced in both directions.

  Bernoulli would make a real difference in a relation where one side is
  much more connected than the other (one-to-many or many-to-one). In that
  case $\text{tph} \gg \text{hpt}$ and Bernoulli would push most corruptions
  towards the tail, generating harder and more informative negatives.

  In our case, since \textit{cites} is close to symmetric, both strategies
  behave the same way and Bernoulli gives no real advantage.

% --------------------------------------------------------------------
\subsection*{A.5 \hspace{0.3em} Disconnected Paper}
% --------------------------------------------------------------------
% TODO
 KGE models learn by minimising a loss computed over training triples.
  For TransE, the loss for a single triple is:
  \[
    \mathcal{L} = |\mathbf{e}_{head} + \mathbf{e}_{rel} - \mathbf{e}_{tail}|
  \]
  Only the entity embeddings that appear in a triple contribute to this
  loss and get updated during training. An isolated paper, one with no
  outgoing or incoming citations, never appears as head or tail in any
  triple. As a result, its loss is never computed, no gradient is ever
  produced for it, and its 128-dimensional embedding vector stays exactly
  at its random initialisation throughout the entire training process.

  While all other papers are gradually moved to meaningful positions in
  the embedding space (papers with similar citation patterns end up close
  to each other), the isolated paper remains at a random point with no
  structural information encoded in it.

  When PCA is applied to project all embeddings to 2D, the principal
  components capture the directions of maximum variance in the trained
  embeddings. Since the isolated paper's embedding is random and has no
  alignment with these structured directions, it projects near the origin
  of the PCA plot,appearing as a lone point close to the centre, far
  from any cluster, with no meaningful neighbours.
% ====================================================================
\section*{B \hspace{0.5em} Exploiting the Global Graph Structure}
% ====================================================================

\subsection*{B.1 \hspace{0.3em} Let's Forget About the Graph}
% TODO

\subsection*{B.2 \hspace{0.3em} Exploiting the Graph Structure}
% TODO

\subsection*{B.3 \hspace{0.3em} Tricking the GNN}
% TODO

\subsection*{B.4 \hspace{0.3em} The More the Merrier?}
% TODO

\subsection*{B.5 \hspace{0.3em} What if We Do Not Have Initial Information?}
% TODO

\subsection*{B.6 \hspace{0.3em} Link Prediction with GNNs}
% TODO

\subsection*{B.7 \hspace{0.3em} Where Are the Embeddings?}
% TODO

\subsection*{B.8 \hspace{0.3em} Embedding Spaces}
% TODO

\subsection*{B.9 \hspace{0.3em} Generating Embeddings for New Papers}
% TODO

\end{document}

    

\end{document}
